%% kapa: exam=1 round=31 year=2020 period=2 form=A subject=appraisal_law subject_ko=감정평가관계법규 seq=23 number=23
%% source: Q-Net 감정평가사 기출문제 (cst00309) artlSeq=5207219 file=extracted/1008/5207219/A형/A형/2020년도 제31회 감정평가사 1차 시험_2교시_A형.pdf
%% conversion: kordoc 4.24.51 -> tools/gen1.py (auto); review=proofread-sonnet; flags=none
\begin{question}{23}
국유재산법령상 국유재산에 관한 설명으로 옳은 것은?
\begin{choices}
  \item 행정재산은 「민법」에 따른 시효취득의 대상이 된다.
  \item 판결에 따라 취득하는 경우에도 사권(私權)이 소멸되지 않은 재산은 국유재산으로 취득하지 못한다.
  \item 총괄청은 일반재산을 보존용재산으로 전환하여 관리할 수 있다.
  \item 직접 공공용으로 사용하기 위하여 국유재산을 관리전환하는 경우에는 유상으로 하여야 한다.
  \item 중앙관서의 장이 국유재산으로 취득한 소유자 없는 부동산은 등기일부터 20년간은 처분할 수 없다.
\end{choices}
\end{question}
\answer{3}
